\chapter{Introduction}
\label{ch:introduction}

\section{The question}
\label{sec:the-question}

What is a Calabi--Yau quantum chiral algebra?

A Calabi--Yau category $\cC$ of dimension $d$ carries a canonical cyclic $\Ainf$-structure: a non-degenerate trace
\[
  \Tr \colon \HH_\bullet(\cC) \longrightarrow k[-d]
\]
on its Hochschild homology, encoding the CY condition as a $d$-dimensional Frobenius structure at the chain level.  The cyclic bar complex $\mathrm{CC}_\bullet(\cC)$ is the primary invariant; it carries an $\Sd$-action (from the $d$-sphere framing of the trace), and its $\Sd$-equivariant structure governs higher-genus amplitudes.

A chiral algebra $A$ on a curve $X$ carries a bar complex $B(A)$, which is a factorization coalgebra on $\Ran(X)$.  The bar differential encodes holomorphic OPE residues; the coproduct encodes topological interval-cutting.  Together they assemble into a Swiss-cheese algebra structure (Volume~II).  At genus $g \geq 1$, the bar complex acquires curvature $\kappa(A) \cdot \omega_g$ from the Hodge bundle, and the full modular structure is controlled by the universal Maurer--Cartan element $\Theta_A$ (Volume~I).

The programme of this work is to construct a precise bridge:
\begin{center}
\begin{tikzcd}[column sep=huge]
  \left\{\begin{array}{c}\text{CY categories with}\\\text{appropriate $E_n$-structure}\end{array}\right\}
  \arrow[r, "\Phi"]
  &
  \left\{\begin{array}{c}\text{quantum chiral algebras}\\\text{with modular trace}\end{array}\right\}
\end{tikzcd}
\end{center}
The functor $\Phi$ should:
\begin{enumerate}[label=(\roman*)]
  \item take a CY category $\cC$ (Fukaya, derived, matrix factorization, or more general) as input;
  \item extract the $E_2$-monoidal structure on $\cC$ (braided tensor product, quantum group action);
  \item produce a chiral algebra $A_\cC$ whose bar complex $B(A_\cC)$ encodes the CY cyclic homology;
  \item realize the CY trace as the modular characteristic $\kappa(A_\cC)$.
\end{enumerate}

\section{The $E_1$/$E_2$ chiral hierarchy}
\label{sec:e1-e2-hierarchy}

The key structural ingredient is the $E_n$-chiral hierarchy:

\begin{itemize}
  \item $E_1$-chiral algebras (Volume~II, Part~VII): associative factorization on $\mathbb{C} \times \mathbb{R}$, encoding the ordered/topological sector.  The representation category of an $E_1$-chiral algebra is monoidal.
  \item $E_2$-chiral algebras (this work): braided factorization on $\mathbb{C} \times \mathbb{C}$, encoding the full holomorphic-holomorphic sector.  The representation category is \emph{braided} monoidal --- the natural habitat of quantum groups.
  \item The passage $E_1 \to E_2$ via Dunn additivity: $E_2 \simeq E_1 \otimes E_1$.  An $E_2$-chiral algebra is an $E_1$-algebra in $E_1$-algebras.
\end{itemize}

The CY condition enters through Kontsevich's identification: a $d$-dimensional CY structure on $\cC$ determines an $\mathbb{S}^d$-framing of the Hochschild complex, hence an $\mathbb{S}^1$-action on $\HH_\bullet(\cC)$.  For $d = 2$ (CY surfaces), this $\mathbb{S}^1$-action is exactly the data of an $E_2$-algebra structure on the cyclic homology --- the braiding.

\section{Relation to Volumes I and II}
\label{sec:relation-vols}

This work depends on and extends the two preceding volumes:

\begin{center}
\renewcommand{\arraystretch}{1.4}
\begin{tabular}{lll}
  \toprule
  \textbf{Volume} & \textbf{Provides} & \textbf{Used here} \\
  \midrule
  I (Modular Koszul Duality) & Bar-cobar machine, $\Theta_A$, $\kappa(A)$ & CY bar complex, modular trace \\
  II (3D HT QFT) & $\SC^{\mathrm{ch,top}}$, PVA descent, DK bridge & $E_1$ sector, braided structure \\
  III (this work) & CY $\to$ chiral functor, $E_2$ theory & --- \\
  \bottomrule
\end{tabular}
\end{center}

\section{CY3 combinatorics as generalized root data}
\label{sec:cy3-root-data}

The deep structural insight animating this work is that the combinatorics of a Calabi--Yau threefold $X$ --- its lattice, intersection form, BPS spectrum --- constitute the \emph{generalized root datum} of what we call a \emph{quantum vertex chiral group} $G(X)$.

Given a CY3 $X$, the generalized root datum $\mathcal{R}(X)$ consists of:
\begin{enumerate}[label=(\roman*)]
  \item A lattice $\Lambda(X)$ extracted from $K(X)$ or $H^*(X)$ with its Euler form;
  \item Real simple roots $\Delta^{\mathrm{re}}$ from the geometry (hyperbolic sublattice for $K3 \times E$; toric diagram for toric CY3);
  \item Imaginary roots $\Delta^{\mathrm{im}} = \Delta^{\mathrm{im}}_0 \sqcup \Delta^{\mathrm{im}}_1$ (even and odd) with multiplicities $\mathrm{mult}(\alpha)$ from BPS counting --- Donaldson--Thomas invariants for $X$, or equivalently the Fourier coefficients of the appropriate automorphic form;
  \item A Weyl group $W(X)$ acting on the positive cone of $\Lambda(X)$;
  \item A Weyl vector $\rho \in \Lambda(X) \otimes \mathbb{Q}$.
\end{enumerate}

Two families of examples illuminate the construction:

\emph{The $K3 \times E$ tower.}  For $X = (S \times E)/(\mathbb{Z}/N\mathbb{Z})$ with $S$ a K3 surface and $E$ an elliptic curve, the lattice is $\Lambda^{3,2} \simeq \Lambda^{1,1} \oplus \Lambda^{1,1} \oplus [2]$ of signature $(3,2)$.  The hyperbolic sublattice $\Lambda^{2,1}_{II}$ with Gram matrix $\bigl(\begin{smallmatrix} 2 & -2 & -2 \\ -2 & 2 & -2 \\ -2 & -2 & 2 \end{smallmatrix}\bigr)$ provides the real roots.  The root multiplicities are the Fourier coefficients $f(n,l)$ of the weak Jacobi form $\phi_{0,1}$ --- the K3 elliptic genus.  The resulting generalized Borcherds--Kac--Moody superalgebra $\mathfrak{g}_{\Delta_5}$ has the Igusa cusp form $\Delta_5$ as its denominator identity.

\emph{Toric CY3 and critical CoHAs.}  For a toric CY3 determined by a fan $\Sigma$, the toric diagram determines a quiver $Q_X$.  The critical cohomological Hall algebra $\mathrm{CoHA}(Q_X)$ is the positive half of an affine super Yangian $Y(\widehat{\mathfrak{g}}_{Q_X})$.  For $\mathbb{C}^3$: the Jordan quiver gives $Y(\widehat{\mathfrak{gl}}_1) \simeq \mathcal{W}_{1+\infty}$.  The toric fan IS the Dynkin data.

\section{Automorphic correction as shadow tower}
\label{sec:automorphic-shadow}

The passage from the ``naive'' Kac--Moody algebra $\mathfrak{g}$ (built from the real root Gram matrix alone) to the full generalized BKM superalgebra $\mathfrak{g}_X$ (with all imaginary roots and their multiplicities) is the \emph{automorphic correction}.  In the language of Volume~I, this passage is precisely the \emph{shadow Postnikov tower}:
\begin{center}
\renewcommand{\arraystretch}{1.4}
\begin{tabular}{lll}
  \toprule
  \textbf{Shadow tower (Vol~I)} & \textbf{BKM language} & \textbf{CY3 geometry} \\
  \midrule
  $\kappa$ (arity 2) & Weyl vector $\rho$ contribution & Leading DT invariant \\
  $C$ (arity 3, cubic) & First imaginary root corrections & Genus-0 3-point function \\
  $Q$ (arity 4, quartic) & Higher imaginary roots & Genus-0 4-point function \\
  Full $\Theta_A$ & Complete automorphic correction & Full BPS generating function \\
  \bottomrule
\end{tabular}
\end{center}
The denominator identity of $\mathfrak{g}_X$ --- which IS the automorphic form ($\Delta_5$ for $K3 \times E$, the DT partition function for toric CY3) --- is the bar-complex Euler product of the quantum vertex chiral group.  The Weyl--Kac--Borcherds character formula is the factorization structure of $B(A_X)$.

\section{Main results}
\label{sec:main-results}

The main theorems of this work are:

\begin{itemize}
  \item \textbf{Theorem CY-A} (CY-to-chiral functor): Construction of $\Phi \colon \CY_d\text{-}\Cat \to E_2\text{-}\mathrm{ChirAlg}$ from CY categories of dimension $d$ to $E_2$-chiral algebras, via the factorization envelope and $\mathbb{S}^d$-framing.  For $d = 3$, $\Phi$ produces a quantum vertex chiral group $G(X)$ whose generalized root datum is $\mathcal{R}(X)$.
  \item \textbf{Theorem CY-B} ($E_2$-chiral Koszul duality): Bar-cobar adjunction in the $E_2$-chiral setting.  The automorphic correction of the BKM superalgebra $\mathfrak{g}_X$ is identified with the shadow tower $\Theta_{A_X}$; the denominator identity equals the bar-complex Euler product.
  \item \textbf{Theorem CY-C} (Quantum group realization): For toric CY3, $\Rep^{E_2}(G(X))$ is braided monoidal equivalent to $\Rep(Y(\widehat{\mathfrak{g}}_{Q_X}))$; for $K3 \times E$, the Sp$_4(\mathbb{Z})$-module structure on the denominator identity recovers the braided structure.  The critical CoHA is the $E_1$-sector (positive half).
  \item \textbf{Theorem CY-D} (Modular CY characteristic): $\kappa(G(X))$ equals the weight of the automorphic form --- the CY Euler characteristic.  The genus-$g$ obstruction $\mathrm{obs}_g = \kappa \cdot \lambda_g$ recovers the genus-$g$ DT/GW invariant on the uniform-weight lane.
\end{itemize}

\section{Guide for the reader}
\label{sec:guide}

Part~I establishes the categorical foundations: CY categories, cyclic $\Ainf$-structures, and the Hochschild calculus.  Part~II develops the $E_1$ and $E_2$ chiral theories.  Part~III constructs the bridge functor $\Phi$ and proves Theorems CY-A through CY-D, with the identification of automorphic correction and shadow tower as the central result.  Part~IV studies quantum groups and their braided factorization structure, including the Drinfeld center as the bulk algebra.  Part~V works through the standard landscape: Fukaya categories, derived categories, matrix factorizations (ADE singularities and W-algebras), and quantum group representations (Kazhdan--Lusztig, affine Yangians, critical CoHAs).  Part~VI connects back to Volumes~I and~II and to the geometric Langlands programme.
